\documentclass[10pt,a4paper]{amsart}
\usepackage{amsmath,amssymb,amsthm,graphicx}
\usepackage{verbatim}
\newenvironment{reference}{\comment}{\endcomment}
%\newenvironment{reference}{}{}
\newenvironment{slogan}{\comment}{\endcomment}
\newenvironment{history}{\comment}{\endcomment}

% For commutative diagrams we use Xy-pic
\usepackage[all]{xy}

% We use 2cell for 2-commutative diagrams.
\xyoption{2cell}
% Theorem environments.
%
\theoremstyle{plain}
\newtheorem{theorem}[subsection]{Theorem}
\newtheorem{proposition}[subsection]{Proposition}
\newtheorem{lemma}[subsection]{Lemma}

\theoremstyle{definition}
\newtheorem{definition}[subsection]{Definition}
\newtheorem{example}[subsection]{Example}
\newtheorem{exercise}[subsection]{Exercise}
\newtheorem{situation}[subsection]{Situation}

\theoremstyle{remark}
\newtheorem{remark}[subsection]{Remark}
\newtheorem{remarks}[subsection]{Remarks}

\numberwithin{equation}{subsection}

% Macros
%
\def\lim{\mathop{\mathrm{lim}}\nolimits}
\def\colim{\mathop{\mathrm{colim}}\nolimits}
\def\Spec{\mathop{\mathrm{Spec}}}
\def\Hom{\mathop{\mathrm{Hom}}\nolimits}
\def\Ext{\mathop{\mathrm{Ext}}\nolimits}
\def\SheafHom{\mathop{\mathcal{H}\!\mathit{om}}\nolimits}
\def\SheafExt{\mathop{\mathcal{E}\!\mathit{xt}}\nolimits}
\def\Sch{\mathit{Sch}}
\def\Mor{\mathop{\mathrm{Mor}}\nolimits}
\def\Ob{\mathop{\mathrm{Ob}}\nolimits}
\def\Sh{\mathop{\mathit{Sh}}\nolimits}
\def\NL{\mathop{N\!L}\nolimits}
\def\CH{\mathop{\mathrm{CH}}\nolimits}
\def\proetale{{pro\text{-}\acute{e}tale}}
\def\etale{{\acute{e}tale}}
\def\QCoh{\mathit{QCoh}}
\def\Ker{\mathop{\mathrm{Ker}}}
\def\Im{\mathop{\mathrm{Im}}}
\def\Coker{\mathop{\mathrm{Coker}}}
\def\Coim{\mathop{\mathrm{Coim}}}

% Boxtimes
%
\DeclareMathSymbol{\boxtimes}{\mathbin}{AMSa}{"02}

%
% Macros for moduli stacks/spaces
%
\def\QCohstack{\mathcal{QC}\!\mathit{oh}}
\def\Cohstack{\mathcal{C}\!\mathit{oh}}
\def\Spacesstack{\mathcal{S}\!\mathit{paces}}
\def\Quotfunctor{\mathrm{Quot}}
\def\Hilbfunctor{\mathrm{Hilb}}
\def\Curvesstack{\mathcal{C}\!\mathit{urves}}
\def\Polarizedstack{\mathcal{P}\!\mathit{olarized}}
\def\Complexesstack{\mathcal{C}\!\mathit{omplexes}}
% \Pic is the operator that assigns to X its picard group, usage \Pic(X)
% \Picardstack_{X/B} denotes the Picard stack of X over B
% \Picardfunctor_{X/B} denotes the Picard functor of X over B
\def\Pic{\mathop{\mathrm{Pic}}\nolimits}
\def\Picardstack{\mathcal{P}\!\mathit{ic}}
\def\Picardfunctor{\mathrm{Pic}}
\def\Deformationcategory{\mathcal{D}\!\mathit{ef}}
\usepackage{cleveref}
\crefname{originalreference}{original reference}{original references}
\Crefname{originalreference}{Original reference}{Original references}
\makeatletter
\newlabel{spaces-morphisms-lemma-base-change-separated}{{03KL}{source}{}{sourcepointer-1}{}}
\newlabel{spaces-morphisms-lemma-base-change-separated@cref}{{[originalreference][0][]03KL}{[1][0][]source}{}{}{}}
\makeatother
\crefname{originaldefinition}{original definition}{original definitions}
\Crefname{originaldefinition}{Original definition}{Original definitions}
\makeatletter
\newlabel{definition-property-morphisms-local}{{03YH}{source}{}{sourcepointer-2}{}}
\newlabel{definition-property-morphisms-local@cref}{{[originaldefinition][0][]03YH}{[1][0][]source}{}{}{}}
\makeatother
\crefname{originalreference}{original reference}{original references}
\Crefname{originalreference}{Original reference}{Original references}
\makeatletter
\newlabel{spaces-topologies-lemma-fppf}{{03Y9}{source}{}{sourcepointer-3}{}}
\newlabel{spaces-topologies-lemma-fppf@cref}{{[originalreference][0][]03Y9}{[1][0][]source}{}{}{}}
\makeatother
\crefname{originallemma}{original lemma}{original lemmas}
\Crefname{originallemma}{Original lemma}{Original lemmas}
\makeatletter
\newlabel{lemma-descending-fppf-property-immersion}{{042U}{source}{}{sourcepointer-4}{}}
\newlabel{lemma-descending-fppf-property-immersion@cref}{{[originallemma][0][]042U}{[1][0][]source}{}{}{}}
\makeatother
\title{Being locally separated is fppf local on the base}
\author{The Stacks Project Authors}
\date{stacks-project:spaces-descent.tex:local-snapshot-20261012; GNU FDL 1.2 or later}
\begin{document}
\maketitle
\noindent\textit{Editorial note.} All delivered mathematical text is the author's own, copied verbatim from the pinned Stacks Project chapter file spaces-descent.tex; no statement or proof step was rewritten, reordered or removed.\par\medskip
\noindent\textit{Editorial note.} The counted span cites one lemma from Morphisms of Spaces, one from Topologies on Spaces, one same-chapter lemma and one definition, recorded in assembly.external\_references with canonical tags as navigation pointers.\par\medskip
\noindent\textit{Allowed standard prerequisites (editorial).}\begin{itemize}
\item Base change of locally separated morphisms. Stacks Project, Morphisms of Spaces, Lemma spaces-morphisms-lemma-base-change-separated. Supplies the direct implication.
\item fppf coverings are stable under base change. Stacks Project, Topologies on Spaces, Lemma spaces-topologies-lemma-fppf. Supplies that \{Y\_i \textbackslash{}times\_Y (X \textbackslash{}times\_Y X) \textbackslash{}to X \textbackslash{}times\_Y X\} is an fppf covering.
\item fppf descent of the immersion property (diagonal form). Stacks Project, Descent on Spaces, Lemma lemma-descending-fppf-property-immersion. Supplies the descent step from the \textbackslash{}Delta\_i to \textbackslash{}Delta.
\end{itemize}\medskip
\ifdefined\hypertarget\hypertarget{sourcepointer-1}{}\fi
\noindent\textit{Source reference.} Original reference 03KL: sources/stacks-chow/tags.txt, source lines 4681--4681. Cross-chapter standard prerequisite cited by the author inside the counted span; it is used as a pointer and its own proof is not part of this item.\par\medskip
\ifdefined\hypertarget\hypertarget{sourcepointer-2}{}\fi
\noindent\textit{Source reference.} Original definition 03YH: sources/stacks-chow/tags.txt, source lines 5184--5184. Cross-chapter standard prerequisite cited by the author inside the counted span; it is used as a pointer and its own proof is not part of this item.\par\medskip
\ifdefined\hypertarget\hypertarget{sourcepointer-3}{}\fi
\noindent\textit{Source reference.} Original reference 03Y9: sources/stacks-chow/tags.txt, source lines 5176--5176. Cross-chapter standard prerequisite cited by the author inside the counted span; it is used as a pointer and its own proof is not part of this item.\par\medskip
\ifdefined\hypertarget\hypertarget{sourcepointer-4}{}\fi
\noindent\textit{Source reference.} Original lemma 042U: sources/stacks-chow/tags.txt, source lines 5341--5341. Cross-chapter standard prerequisite cited by the author inside the counted span; it is used as a pointer and its own proof is not part of this item.\par\medskip
\begin{lemma}
\label{lemma-descending-fppf-property-locally-separated}
The property $\mathcal{P}(f) =$``$f$ is locally separated''
is fppf local on the base.
\end{lemma}

\begin{proof}
A base change of a locally separated morphism is locally separated, see
Morphisms of Spaces,
Lemma \ref{spaces-morphisms-lemma-base-change-separated}.
Hence the direct implication in
Definition \ref{definition-property-morphisms-local}.

\medskip\noindent
Let $\{Y_i \to Y\}_{i \in I}$ be an fppf covering of algebraic spaces over $S$.
Let $f : X \to Y$ be a morphism of algebraic spaces over $S$.
Assume each base change $X_i := Y_i \times_Y X \to Y_i$ is locally separated.
This means that each of the morphisms
$$
\Delta_i :
X_i
\longrightarrow
X_i \times_{Y_i} X_i = Y_i \times_Y (X \times_Y X)
$$
is an immersion. The base change of a fppf covering is an
fppf covering, see
Topologies on Spaces, Lemma \ref{spaces-topologies-lemma-fppf}
hence $\{Y_i \times_Y (X \times_Y X) \to X \times_Y X\}$
is an fppf covering of algebraic spaces. Moreover, each
$\Delta_i$ is the base change of the morphism
$\Delta : X \to X \times_Y X$. Hence it follows from
Lemma \ref{lemma-descending-fppf-property-immersion}
that $\Delta$ is an immersion, i.e., $f$ is locally separated.
\end{proof}

\end{document}
